\chapter{Complexity — Input Device}
%%% generated by the blueprint pipeline from TCSlib.Complexity.TuringMachine.Build.InputDevice
\section{Overview}

Declarations of the Lean module \texttt{TCSlib.Complexity.TuringMachine.Build.InputDevice}.

\section{Declarations}

\begin{definition}[Input symbol at a position]
\lean{Turing.inputCell}
\label{Turing.inputCell}
\leanok
For a list $y$ of symbols and a natural number $i$, define the input cell at position
$i$ to be blank when $i=0$, and otherwise to be the symbol of $y$ at zero-based index
$i-1$, if that index exists.
\end{definition}

\begin{theorem}[Input Symbol as an Input Cell]
\lean{Turing.Cfg.inputSymbol_eq_inputCell}
\label{Turing.Cfg.inputSymbol_eq_inputCell}
\leanok
\uses{Turing.Cfg, Turing.Cfg.inputSymbol, Turing.inputCell, Turing.inputSymbolInner}
\difficulty{5}
For any configuration $c$ with input word $y$, the optional symbol scanned by its input
head equals the content of input cell $c$'s head position, where input cell $0$ is the
left boundary blank and the input symbols occupy the subsequent positions.
\end{theorem}

\begin{definition}[Input head at a position]
\lean{Turing.NativeAt}
\label{Turing.NativeAt}
\leanok
\uses{Turing.Cfg}
For a configuration of a machine with $K$ work tapes on an input list $y$, and a
position $i$ among the $|y|+2$ input-head positions, the configuration is at $i$ when
its input-head position equals $i$.
\end{definition}

\begin{theorem}[Scanned Symbol at a Given Position]
\lean{Turing.NativeAt.inputSymbol_eq}
\label{Turing.NativeAt.inputSymbol_eq}
\leanok
\uses{Turing.Cfg, Turing.Cfg.inputSymbol, Turing.Cfg.inputSymbol_eq_inputCell, Turing.NativeAt, Turing.inputCell}
\difficulty{1}
Let $d$ be a configuration on input word $y$, and let $i$ be one of the $|y|+2$
input-head positions. If the input head of $d$ is at position $i$, then the symbol read
by the head is the input-cell value at position $i$.
\end{theorem}

\begin{theorem}[Head Position After Movement]
\lean{Turing.NativeAt.move}
\label{Turing.NativeAt.move}
\leanok
\uses{Turing.Cfg, Turing.NativeAt, Turing.moveInputPos}
\difficulty{3}
Let $d$ be a configuration on input word $y$, and let $i$ be an input-head position. If
the head of $d$ is at $i$, then after moving the head by any sign $m\in\{-1,0,1\}$ with
boundary clamping, the head is at the correspondingly moved position of $i$.
\end{theorem}

\begin{theorem}[Initial input-head position]
\lean{Turing.nativeAt_init}
\label{Turing.nativeAt_init}
\leanok
\uses{Turing.Cfg, Turing.Cfg.init, Turing.NativeAt}
\difficulty{1}
For every natural number $K$, state set $S$, symbol alphabet $\Gamma$, starting state
$q_0\in S$, and input word $y\in\Gamma^*$, the initial configuration on $y$ has its
input head at position $1$.
\end{theorem}

\begin{theorem}[When a shifted input cell is blank]
\lean{Turing.inputCell_sub_eq_none_iff}
\label{Turing.inputCell_sub_eq_none_iff}
\leanok
\uses{Turing.inputCell}
\difficulty{4}
Let $z$ be a Boolean word and let $\ell,p$ be natural numbers with $0<\ell\le p\le
|z|+1$. The input cell of $z$ at index $p-\ell$ is blank exactly when $p=\ell$.
\end{theorem}

\begin{theorem}[Input cells after dropping a prefix]
\lean{Turing.inputCell_of_drop}
\label{Turing.inputCell_of_drop}
\leanok
\uses{Turing.inputCell}
\difficulty{3}
Let $z$ and $y$ be Boolean words, let $\ell$ be a natural number, and suppose deleting
the first $\ell$ symbols of $z$ gives $y$. For every nonzero index $i$, the symbol read
at index $i$ in $y$ equals the symbol read at index $\ell+i$ in $z$, with indices
outside a word yielding no symbol.
\end{theorem}

\begin{definition}[Ruler tape configuration]
\lean{Turing.RulerAt}
\label{Turing.RulerAt}
\leanok
\uses{Turing.Cfg, Turing.FinTM.bufferTape}
Let $K$ be a natural number, let $y$ and $z$ be finite Boolean strings, let $d$ be a
configuration of a $K$-tape machine on input $z$, let $r$ be one of its work-tape
indices, let $i$ be an input-head position among the $|y|+2$ positions, and let $\ell$ be a
natural number. The ruler condition holds if $0<\le |z|$, the string obtained from $z$
by deleting its first $\ell$ symbols is $y$, and the input head is at position $\ell+i$. In
addition, for every integer $p$ with $-1\le p\le \ell$, the selected work tape at offset
$p$ from its head agrees with the tape encoding a string of $\ell$ true symbols.
\end{definition}

\begin{theorem}[Length of a nested pair encoding]
\lean{Turing.pairEncode_pairEncode_length}
\label{Turing.pairEncode_pairEncode_length}
\leanok
\uses{Turing.pairEncode, Turing.universal_pair_length}
\difficulty{3}
For any bit strings $u$, $\alpha$, and $x$, the length of the encoding of
$(u,\langle\alpha,x\rangle)$ is $2|u|+2|\alpha|+4+|x|$, where $|w|$ denotes the length
of a bit string $w$.
\end{theorem}

\begin{theorem}[Recovering a nested pair encoding]
\lean{Turing.pairEncode_pairEncode_drop}
\label{Turing.pairEncode_pairEncode_drop}
\leanok
\uses{Turing.pairEncode, Turing.universal_pair_length}
\difficulty{4}
For any bit strings $u$, $\alpha$, and $x$, removing the first $2|u|+2|\alpha|+4$ bits
from the encoding of $(u,\langle\alpha,x\rangle)$ leaves exactly $x$, where $|w|$
denotes the length of a bit string $w$.
\end{theorem}

\begin{theorem}[Ruler-boundary test machine]
\lean{Turing.exists_rulerTestTM}
\label{Turing.exists_rulerTestTM}
\leanok
\uses{Turing.Cfg, Turing.FinTM.bufferTape, Turing.MultiTapeTM, Turing.MultiTapeTM.runFrom}
There exist a finite state set $P$, a deterministic one-work-tape binary Turing machine
$R$, entry and distinct yes/no states, and a constant $C$. For every binary input $z$,
ruler length $\ell>0$ not exceeding the input-head position, and configuration $d$ in
the entry state whose work tape contains a ruler of $\ell$ ones at offsets $-1$ through
$\ell$, the machine reaches within $C(\ell+1)$ steps the configuration identical to $d$
except that its state is yes exactly when the input-head position equals $\ell$, and no
otherwise. It visits neither yes nor no earlier, and its work-head displacement at every
time up to arrival has absolute value at most $\ell+1$.
\end{theorem}

\begin{theorem}[Existence of a Ruler Device]
\lean{Turing.exists_rulerDeviceTM}
\label{Turing.exists_rulerDeviceTM}
\leanok
\uses{Turing.Cfg, Turing.MultiTapeTM, Turing.MultiTapeTM.runFrom, Turing.RulerAt, Turing.inputCell, Turing.moveInputPos}
There exist a finite state set $P$ with decidable equality, a one-work-tape Boolean
Turing machine $R$ with state set $P$, states $readEntry,moveExit\in P$, an injective
map $readExit:\mathrm{Option}(\{\mathrm{false},\mathrm{true}\})\to P$ whose image
excludes $moveExit$, a map $moveEntry:\{-1,0,1\}\to P$, and a constant $C\in\mathbb N$
with the following property. For any Boolean words $y,z$, any $\ell$ with $0<\ell\le
|z|$, any configuration $d$ of $R$ on input $z$, and any input-head position $i$ among
the $|y|+2$ positions, suppose deleting the first $\ell$ symbols of $z$ gives $y$, the
input head of $d$ is at $\ell+i$, and its work tape at every offset $p$ from $-1$
through $\ell$ from its head contains the corresponding cell of the tape encoding $\ell$
true symbols. If $d$ is in state $readEntry$, then within $C(\ell+1)$ steps the machine
reaches the otherwise unchanged configuration with state $readExit(inputCell(y,i))$;
before that time it does not enter any $readExit$ state, and throughout the run its work
head stays within distance $\ell+1$ of its initial position. For every movement sign
$m$, if $d$ is in state $moveEntry(m)$, then within $C(\ell+1)$ steps it reaches state
$moveExit$, with input head at $\ell+moveInputPos(i,m)$, unchanged work tape, work-head
position, and output; before that time it does not enter $moveExit$, and throughout the
run its work head stays within distance $\ell+1$ of its initial position.
\end{theorem}

\begin{definition}[Configuration with an added work tape]
\lean{Turing.rulerHostCfg}
\label{Turing.rulerHostCfg}
\leanok
\uses{Turing.Cfg}
Let $m$ and $\ell$ be natural numbers, let $S$ and $Q$ be types, and let $y$ and $z$ be
Boolean lists with $|z|=\ell+|y|$. Given a state $q\in\mathrm{Option}(Q)$, a
configuration $c$ with $m$ work tapes on input $y$, a tape
$r:\mathbb{Z}\to\mathrm{Option}(\mathrm{Bool})$, and a head position $p\in\mathbb{Z}$,
form a configuration with $m+1$ work tapes on input $z$: its state is $q$, its
input-head position is shifted by $\ell$, its work tapes are those of $c$ followed by
$r$, its work-head positions are those of $c$ followed by $p$, and its output is the
output of $c$.
\end{definition}

\begin{theorem}[Ruler-host simulation]
\lean{Turing.exists_rulerHostTM}
\label{Turing.exists_rulerHostTM}
\leanok
\uses{Turing.Cfg, Turing.FinTM.bufferTape, Turing.MultiTapeTM, Turing.MultiTapeTM.runFrom, Turing.MultiTapeTM.visitedByTapeHead, Turing.rulerHostCfg}
For every deterministic $m$-work-tape binary Turing machine $M$ with a finite state set
$S$ and decidable state equality, there are a finite state set $Q$, a deterministic
binary machine $H$ with $m+1$ work tapes, an injection $\mathrm{inj}:S\times\{0,1\}\to
Q$, and a positive integer $C$ with the following property. Let $y,z$ be binary strings,
let $\ell>0$ satisfy $|z|=\ell+|y|$ and $z$'s suffix after its first $\ell$ symbols
equal $y$, and let the extra tape contain a ruler of $\ell$ ones from offset $-1$
through $\ell$, with head position $rp$. For any configuration $c$ of $M$ on input $y$
and bit $b$ indicating whether its input head is at position zero, and any time $t$, the
host machine started on the paired input and configuration formed by shifting the
input-head position by $\ell$, preserving the original tapes and output, and placing the
ruler on its extra tape, reaches after some $t'$ with $t\le t'\le C(\ell+1)t$ the
corresponding configuration of $M$ after $t$ steps, with the state encoded using a bit
indicating whether that simulated input head is at zero. If $M$ has not halted before
$t$, the host has not halted before $t'$; at every intermediate host time $u\le t'$, its
original tapes and output agree with some simulated configuration at a time $v\le t$,
the ruler tape is unchanged, and its ruler head remains in $[rp-1,rp+\ell]$. By time
$t'$, each original tape has visited exactly the cells visited by the corresponding tape
of $M$ through time $t$, and the ruler head has visited only cells in that same
interval.
\end{theorem}
